\documentclass[11pt,a4paper]{article}

\usepackage[utf8]{inputenc}
\usepackage[T1]{fontenc}
\usepackage{amsmath,amssymb,amsthm,mathtools}
\usepackage{bbm}
\usepackage{bm}
\usepackage{hyperref}
\usepackage[margin=2.5cm]{geometry}
\usepackage{enumitem}
\usepackage{booktabs}
\usepackage{natbib}
\bibliographystyle{abbrvnat}

\newtheorem{theorem}{Theorem}[section]
\newtheorem{proposition}[theorem]{Proposition}
\newtheorem{lemma}[theorem]{Lemma}
\newtheorem{corollary}[theorem]{Corollary}
\theoremstyle{definition}
\newtheorem{definition}[theorem]{Definition}
\newtheorem{assumption}{Assumption}
\theoremstyle{remark}
\newtheorem{remark}[theorem]{Remark}

\newcommand{\R}{\mathbb{R}}
\newcommand{\E}{\mathbb{E}}
\newcommand{\cP}{\mathcal{P}}
\newcommand{\cS}{\mathcal{S}}
\newcommand{\vtheta}{\boldsymbol{\theta}}
\newcommand{\vw}{\mathbf{w}}
\newcommand{\vx}{\mathbf{x}}
\newcommand{\va}{\mathbf{a}}
\newcommand{\vp}{\mathbf{p}}
\newcommand{\ve}{\mathbf{e}}
\newcommand{\vone}{\mathbf{1}}
\newcommand{\vs}{\mathbf{s}}
\newcommand{\vz}{\mathbf{z}}
\newcommand{\simplex}{\Delta^{p-1}}
\newcommand{\entmax}{\operatorname{entmax}}
\newcommand{\supp}{\operatorname{supp}}
\newcommand{\Var}{\operatorname{Var}}
\DeclareMathOperator*{\argmax}{arg\,max}
\newcommand{\dd}{\mathrm{d}}
\newcommand{\dmu}{\dd\mu}
\newcommand{\norm}[1]{\left\|#1\right\|}
\newcommand{\abs}[1]{\left|#1\right|}

\title{%
  Dynamical Constraints on Feature Selection\\
  in Soft Decision Stump Ensembles
}
\author{
  Ryuichi Kanoh
  \and
  Mahito Sugiyama
}
\date{\today}

\begin{document}
\maketitle

\begin{abstract}
  We study the mean-field dynamics of soft decision stump ensembles
  whose splitting directions are parameterized via entmax, a
  sparsity-inducing map on the probability simplex.
  %
  We identify two competing dynamical constraints that govern
  sparse feature selection in this model.
  %
  \textbf{Irreversibility} (Theorem~\ref{thm:absorption}):
  any feature excluded from the entmax support receives exactly zero
  gradient and its weight freezes; moreover, the entmax threshold
  rises as in-support features strengthen, making re-entry impossible.
  This creates absorbing dead-feature states from which the dynamics
  cannot escape.
  %
  \textbf{Bottleneck} (Theorem~\ref{thm:separation}):
  even when all features are in the support (high temperature),
  the relevance gap between relevant and irrelevant features grows
  only at rate $O(1/(\tau p))$, where $\tau$ is the entmax
  temperature and $p$ is the ambient dimension.
  %
  Together, these constraints imply that sparse feature selection
  in differentiable stumps is caught between a failure mode
  (premature sparsity kills features permanently) and a bottleneck
  (conservative temperature slows selection to $O(\tau^2 p)$ time).
  %
  As a corollary, we derive a quantitative lower bound on the
  temperature needed to avoid dead features, and an explicit
  pruning time for the surviving regime.
  %
  Beyond the early phase, we show that the residual coupling
  between stumps drives a \emph{signal crossing}: the initially
  strongest feature's effective signal decays as it is learned,
  eventually falling below that of weaker features
  (Proposition~\ref{prop:crossing}).
  This dynamic rebalancing is a genuinely collective phenomenon
  that does not occur under independent training.
\end{abstract}

% ============================================================
\section{Introduction}
\label{sec:intro}
% ============================================================

Soft (differentiable) decision trees trained by gradient methods
have gained attention as an alternative to classical greedy
trees~\citep{frosst2017distilling,hazimeh2020tree,popov2020neural}.
A natural approach to imposing axis-aligned feature selection in
such models is to apply a sparsity-inducing map---such as
entmax~\citep{peters2019sparse,blondel2020learning}---to the
splitting weight vector.
As the entmax temperature decreases, the output concentrates on
fewer features, approaching a one-hot (axis-aligned) split.

This paper asks: \emph{what are the dynamical constraints on
feature selection when soft stumps are trained by gradient flow
with entmax gating?}

We work in the mean-field regime
\citep{mei2018mean,chizat2018global}, where the number of stumps
tends to infinity and the ensemble's evolution is described by a
gradient flow on the space of probability measures over stump
parameters.
We focus on depth-one trees (stumps) under a sparse additive
target with independent covariates---the simplest regime in which
the feature-selection dynamics can be rigorously analyzed.

\paragraph{Main findings.}
We identify two dynamical constraints that together determine the
feasibility and speed of sparse feature selection.

\begin{enumerate}[label=(\roman*),leftmargin=*]
  \item \textbf{Irreversible feature exclusion}
        (Theorem~\ref{thm:absorption}).
        The zero-column structure of the entmax Jacobian means that
        any feature outside the support receives exactly zero gradient,
        freezing its weight.
        We show further that the entmax threshold $\lambda^*$ is
        strictly increasing when in-support features have positive
        gradient drift (which they do, by our separation result),
        so excluded features can never re-enter the support.
        This creates \emph{absorbing dead-feature states}.
        %
        As a corollary, if the temperature is too low at initialization
        and a relevant feature starts outside the support, the
        ensemble \emph{provably cannot recover} that feature,
        regardless of training time
        (Corollary~\ref{cor:failure}).

  \item \textbf{Quantitative bottleneck for gap growth}
        (Theorem~\ref{thm:separation}).
        Even in the favorable regime where all features are in the
        support (high temperature), the gradient signal for
        each relevant feature is $O(1)$ while that for each
        irrelevant feature is $O(1/p)$.
        Through the entmax Jacobian, this separation translates into
        a relevance-gap growth rate of only $O(1/(\tau p))$.
        Successful feature selection therefore requires
        $O(\tau^2 p / \kappa_{\min})$ training time---a
        quantitative bottleneck that cannot be circumvented by
        algorithmic tricks within this dynamics.

  \item \textbf{Residual-driven specialization}
        (Proposition~\ref{prop:crossing}).
        Beyond the early phase, the mean-field coupling through
        the shared residual produces \emph{signal crossing}:
        as the ensemble learns the strongest feature, that feature's
        effective signal decays, eventually falling below weaker
        features.
        This rebalancing is a collective phenomenon that does not
        occur under independent stump training, demonstrating that
        the mean-field interaction has a functional role in
        producing feature-level diversity.
\end{enumerate}

Together, (i) and (ii) determine the feasible operating regime
for feature selection: premature sparsity kills features permanently,
while conservative temperature slows selection to $O(\tau^2 p)$
time.
Result~(iii) shows that, beyond these constraints, the shared
residual actively drives specialization across features.

% ============================================================
\section{Problem Setup}
\label{sec:setup}
% ============================================================

\subsection{Data Model}

\begin{assumption}[Sparse Additive Model]
  \label{asm:data}
  \leavevmode
  \begin{enumerate}[label=(\alph*)]
    \item $x_1, \dots, x_p$ are mutually independent with
          $\E[x_j] = 0$, $\E[x_j^2] = 1$,
          $\E[x_j^4] \le \kappa_4$ for all $j$.
    \item $y = \sum_{j \in S} g_j(x_j) + \varepsilon$,
          $S \subset [p]$, $|S| = k$,
          $\E[g_j(x_j)] = 0$, $\|g_j\|_{L^2} \le B_g$,
          $\varepsilon \perp \vx$, $\E[\varepsilon] = 0$.
  \end{enumerate}
\end{assumption}

\begin{assumption}[First-Order Signal]
  \label{asm:signal}
  $\kappa_j \coloneqq \E[g_j(x_j) x_j] \neq 0$ for all $j \in S$.
  Write $\kappa_{\min} = \min_{j \in S} |\kappa_j|$.
\end{assumption}

\subsection{Soft Decision Stump with Entmax}

A stump is parameterized by $\vtheta = (\vw, b, c) \in \R^{p+2}$:
\begin{equation}
  T(\vx;\,\vtheta)
    = c \cdot \sigma_\beta\!\big(
        \va_\alpha^\tau(\vw)^\top \vx - b
      \big),
  \label{eq:stump}
\end{equation}
where $\va_\alpha^\tau(\vw) = \entmax_\alpha(\vw/\tau) \in \simplex$.

\begin{definition}[Entmax]
  For $\alpha > 1$:
  $\entmax_\alpha(\vz) = \argmax_{\vp \in \simplex}
  \{\vp^\top\vz - \Omega_\alpha(\vp)\}$
  with $\Omega_\alpha(\vp) = \frac{1}{\alpha(\alpha-1)}
  \sum_j(p_j^\alpha - p_j)$.
  The solution is
  $[\entmax_\alpha(\vz)]_j
  = [(\alpha{-}1)(z_j - \lambda^*)]_+^{1/(\alpha-1)}$
  where $\lambda^*$ is uniquely determined by
  $\sum_j [(\alpha{-}1)(z_j - \lambda^*)]_+^{1/(\alpha-1)} = 1$.
  The support is $\cS(\vz) = \{j : z_j > \lambda^*\}$.
\end{definition}

\begin{lemma}[Entmax properties {\citep{blondel2020learning}}]
  \label{lem:entmax}
  For $\alpha \in (1,2)$: $\entmax_\alpha$ is differentiable with
  continuous Jacobian and is $1$-Lipschitz.
  The Jacobian satisfies $[J_\alpha(\vz)]_{ij} = 0$ whenever
  $j \notin \cS(\vz)$ (zero columns outside the support).
\end{lemma}

\subsection{Mean-Field Formulation}

\begin{assumption}[Bounded Parameters]
  \label{asm:bounded}
  Parameters lie in
  $\Theta = \{(\vw,b,c): \|\vw\|_\infty \le W,\,|b|\le B,\,|c|\le C\}$.
\end{assumption}

The mean-field predictor is
$f_\mu(\vx) = \int T(\vx;\,\vtheta)\,\dmu(\vtheta)$,
and each particle evolves by
\begin{equation}
  \dot{\vtheta}(t) = \E_{\vx}\!\big[
    r_{\mu_t}(\vx) \cdot \nabla_{\vtheta} T(\vx;\,\vtheta(t))
  \big],
  \qquad
  r_\mu = f^* - f_\mu.
  \label{eq:ode}
\end{equation}

\begin{proposition}[Mean-Field Limit]
  \label{prop:mf}
  Under Assumptions~\ref{asm:data}--\ref{asm:bounded}, the map
  $\vtheta \mapsto T(\vx;\,\vtheta)$ is Lipschitz with constant
  $L_T = 1 + C(\sqrt{p}/\tau+1)/(4\beta)$.
  The drift is bounded on $\Theta$, so for any $T>0$,
  choosing $W,B,C$ large enough keeps the flow in $\Theta$ on
  $[0,T]$.
  The propagation-of-chaos framework of \citet{mei2018mean} gives
  $\sup_{t\in[0,T]} W_2(\mu_t^{(M)},\mu_t) \to 0$ in probability.
\end{proposition}

% ============================================================
\section{Main Results}
\label{sec:main}
% ============================================================

\subsection{Irreversible Feature Exclusion}

This is the failure-mode result.

\begin{theorem}[Dead-Feature Absorption]
  \label{thm:absorption}
  Fix $\alpha \in (1,2)$, $\tau > 0$.
  Consider the mean-field gradient flow~\eqref{eq:ode}.
  If at time $t_0 \ge 0$, a coordinate $j$ satisfies
  $j \notin \cS(\vw(t_0)/\tau)$, then:
  \begin{enumerate}[label=(\alph*)]
    \item \textbf{Weight freezing.}
          $\dot{w}_j(t) = 0$ for all $t \ge t_0$, so
          $w_j(t) = w_j(t_0)$ for all $t \ge t_0$.
    \item \textbf{Threshold monotonicity.}
          If some in-support feature $m \in \cS(\vw(t_0)/\tau)$
          has $\dot{w}_m(t_0) > 0$, then the entmax threshold
          satisfies $\dot{\lambda}^*(t_0) > 0$.
    \item \textbf{Permanent exclusion.}
          If $\dot{w}_m(t) > 0$ for at least one
          $m \in \cS(\vw(t)/\tau)$ for all $t \ge t_0$,
          then $j \notin \cS(\vw(t)/\tau)$ for all $t \ge t_0$.
  \end{enumerate}
\end{theorem}

\begin{proof}
  \emph{Part (a).}
  The $w_j$-component of~\eqref{eq:ode} involves
  $\frac{\partial T}{\partial w_j}
  = c\,\sigma_\beta'(\va^\top\vx - b)
  \sum_l \frac{1}{\tau}[J_\alpha(\vw/\tau)]_{lj}\,x_l$.
  By Lemma~\ref{lem:entmax},
  $[J_\alpha(\vw/\tau)]_{lj} = 0$ for all $l$ when
  $j \notin \cS(\vw/\tau)$.
  Hence $\frac{\partial T}{\partial w_j} = 0$, giving
  $\dot{w}_j = 0$.

  \emph{Part (b).}
  Write $z_m = w_m/\tau$ and
  $\vp = \entmax_\alpha(\vz)$ with $s_m = p_m^{2-\alpha}$.
  Differentiating the normalization
  $F(\lambda^*) \coloneqq
  \sum_{m \in \cS} [(\alpha{-}1)(z_m - \lambda^*)]^{1/(\alpha-1)}
  = 1$
  with respect to time:
  \[
    0 = \frac{\dd F}{\dd t}
      = \sum_{m \in \cS} s_m\,(\dot{z}_m - \dot{\lambda}^*)
      = \sum_{m \in \cS} s_m\,\dot{z}_m
        - (\vone^\top\vs)\,\dot{\lambda}^*.
  \]
  Solving:
  \begin{equation}
    \dot{\lambda}^*
    = \frac{\sum_{m \in \cS} s_m\,\dot{z}_m}{\vone^\top\vs}
    = \frac{\sum_{m \in \cS} s_m\,\dot{w}_m}
           {\tau\,\vone^\top\vs}.
    \label{eq:lambda-dot}
  \end{equation}
  Since $s_m > 0$ for $m \in \cS$ and $\dot{w}_m > 0$
  for at least one such~$m$ (while $\dot{w}_m \ge 0$ for the others
  during the early phase---see Theorem~\ref{thm:separation}(c)),
  we get $\dot{\lambda}^* > 0$.

  \emph{Part (c).}
  By (a), $w_j(t) = w_j(t_0)$ for $t \ge t_0$,
  so $z_j(t) = z_j(t_0)$ is constant.
  By (b), $\lambda^*(t)$ is strictly increasing as long as the
  positivity condition holds.
  Since $j \notin \cS$ at $t_0$ means $z_j(t_0) \le \lambda^*(t_0)$,
  and $\lambda^*$ increases while $z_j$ stays fixed,
  $z_j(t) < \lambda^*(t)$ for all $t > t_0$.
\end{proof}

\begin{corollary}[Low-Temperature Failure]
  \label{cor:failure}
  If the temperature $\tau$ is low enough that some
  $j^* \in S$ satisfies $j^* \notin \cS(\vw(0)/\tau)$,
  and the positivity condition of
  Theorem~\ref{thm:absorption}(c) holds, then:
  \begin{enumerate}[label=(\alph*)]
    \item $j^*$ is permanently excluded from the support.
    \item The predictor $f_{\mu_t}$ cannot represent
          $g_{j^*}(x_{j^*})$.
    \item The excess risk satisfies
          $R[\mu_t] - \inf_\mu R[\mu]
          \ge \frac{1}{2}\|g_{j^*}\|_{L^2}^2$
          for all $t \ge 0$.
  \end{enumerate}
\end{corollary}

\begin{proof}
  Part (a) is Theorem~\ref{thm:absorption}(c).
  For (b): since $j^*$ is excluded, each stump's output
  $T(\vx;\,\vtheta)$ depends on $x_{j^*}$ only through
  $a_{j^*} x_{j^*}$ with $a_{j^*} = 0$ (entmax assigns zero weight
  to excluded features).
  Thus $f_{\mu_t}$ is constant in $x_{j^*}$.
  For (c):
  $\|f^* - f_{\mu_t}\|_{L^2}^2
  \ge \|\E[f^* - f_{\mu_t} \mid x_{j^*}]\|_{L^2}^2
  = \|g_{j^*} - \E[f_{\mu_t} \mid x_{j^*}]\|_{L^2}^2
  = \|g_{j^*}\|_{L^2}^2$,
  since $\E[f_{\mu_t} \mid x_{j^*}]$ is constant (independent of
  $x_{j^*}$) and $\E[g_{j^*}] = 0$.
\end{proof}

\begin{remark}[When does low-temperature failure occur?]
  \label{rem:failure-condition}
  For symmetric initialization $\vw(0) = w_0\vone$,
  $\cS(\vw(0)/\tau) = [p]$ for all $\tau > 0$
  (all features in support), so failure does not occur.
  For random initialization $w_j(0) \overset{iid}{\sim}
  \mathcal{N}(0, \sigma_w^2)$, the spread
  $\max_j w_j(0) - \min_j w_j(0) = \Theta(\sigma_w\sqrt{\log p})$,
  and entmax excludes coordinates whose $w_j/\tau$ falls below
  $\lambda^*$.
  A sufficient condition for full support is
  $\tau \ge C_\alpha\,\sigma_w\sqrt{\log p}$
  for a constant $C_\alpha$ depending on~$\alpha$.
  Below this threshold, features may be permanently lost.
\end{remark}

\subsection{Quantitative Bottleneck for Feature Selection}

This is the rate result.
We use the same notation as before:
$G_j(\vtheta;\,\mu) = \E[r_\mu(\vx)\,\sigma_\beta'(\va^\top\vx-b)\,x_j]$.

\begin{assumption}[High Temperature]
  \label{asm:hightemp}
  $\tau$ is large enough that $\cS(\vw/\tau) = [p]$ for all
  $\|\vw\|_\infty \le W$.
\end{assumption}

\begin{definition}[Early Phase]
  \label{def:early}
  $t_{\mathrm{ep}} = \sup\{t \ge 0 :
  \|f_{\mu_s}\|_{L^2} \le \kappa_{\min}/2
  \;\forall s \in [0,t]\}$.
\end{definition}

\begin{theorem}[Feature Gradient Separation]
  \label{thm:separation}
  Under Assumptions~\ref{asm:data}--\ref{asm:hightemp},
  for a particle with $c > 0$ and for $t \in [0, t_{\mathrm{ep}})$:
  \begin{enumerate}[label=(\alph*)]
    \item \textbf{Relevant features:}
    $G_j = \E[r_j^{(\mu)}(x_j)\,x_j]
    \cdot \E[\sigma_\beta'(Z_j{-}b)] + O(1/p)$
    for $j \in S$,
    where $Z_j = \sum_{l \neq j} a_l x_l \perp x_j$ and
    $\E[r_j^{(\mu)} x_j] \ge \kappa_{\min}/2$ during the early phase.

    \item \textbf{Irrelevant features:}
    $|G_l| \le C_2/p$ for $l \notin S$.

    \item \textbf{Weight dynamics:}
    $\dot{w}_j = \frac{c}{\tau p}\,G_j + O(1/p^2)$ for $j \in S$;
    \quad $\dot{w}_l = O(1/p^2)$ for $l \notin S$.

    \item \textbf{Gap growth:}
    $\frac{\dd}{\dd t}\Delta(\vw)
    = \frac{c}{\tau p}\bar{G}_S + O(1/p^2)$,
    where $\Delta(\vw) = \frac{1}{k}\sum_{j \in S} w_j
    - \frac{1}{p-k}\sum_{l \notin S} w_l$
    and $\bar{G}_S > 0$ during the early phase.
  \end{enumerate}
\end{theorem}

The proof is in Section~\ref{sec:proofs}.
Parts (a) and (b) do not assume that $r_\mu$ decomposes
additively; the proof uses only independence of features and
$\E[x_j] = 0$.

\begin{remark}[The $1/(\tau p)$ rate is the price of exploration]
  The factor $1/p$ in the gap growth rate arises because
  the high-temperature entmax assigns weight $1/p$ to each feature,
  diluting the gradient signal.
  The factor $1/\tau$ comes from the chain rule through
  $\entmax_\alpha(\vw/\tau)$.
  Neither factor can be removed without exiting the
  high-temperature regime (which risks dead features by
  Theorem~\ref{thm:absorption}).
  This is the fundamental bottleneck:
  exploration (high $\tau$, full support) is necessary to avoid
  irreversible feature loss, but it directly slows feature
  selection.
\end{remark}

\subsection{Pruning Time}

Combining the two constraints:

\begin{lemma}[Entmax Support Threshold]
  \label{lem:threshold}
  If $\min_{j \in S} w_j - \max_{l \notin S} w_l
  \ge \Lambda\,\tau$ with
  $\Lambda = \Lambda(\alpha,k) = 1/((\alpha{-}1)k^{\alpha-1})$,
  then $\cS(\vw/\tau) \subseteq S$.
\end{lemma}

\begin{proof}
  Since all $z_j = w_j/\tau$ for $j \in S$ exceed all $z_l$ for
  $l \notin S$ by at least $\Lambda$, $S \subseteq \cS(\vz)$.
  The normalization gives
  $\lambda^* \ge \min_{j \in S} z_j - \Lambda$
  $\ge \max_{l \notin S} z_l$,
  so no $l \notin S$ is in the support.
\end{proof}

\begin{proposition}[Conditional Pruning Time]
  \label{prop:pruning}
  Under symmetric initialization $\vw(0) = w_0\vone$ and
  fixed~$\tau$, conditional on $t^* \le t_{\mathrm{ep}}$:
  the support reduces to $S$ by time
  \begin{equation}
    t^* = O\!\left(\frac{\tau^2 p}{c\,\kappa_{\min}}\right).
    \label{eq:tstar}
  \end{equation}
\end{proposition}

\begin{proof}
  By Theorem~\ref{thm:separation}(c),
  $w_j(t) - w_l(t) \ge \gamma t - O(t/p^2)$ for $j \in S$,
  $l \notin S$, where
  $\gamma = \frac{c\kappa_{\min}}{2\tau p}
  \inf_b \E[\sigma_\beta'(\bar{X}-b)]$.
  Setting $\gamma t^* = \Lambda\tau$ gives $t^*$.
\end{proof}

\subsection{The Tradeoff}

Theorems~\ref{thm:absorption} and~\ref{thm:separation} together
determine the feasible regime.
For random initialization with spread
$\delta_0 = \Theta(\sigma_w\sqrt{\log p})$:
\begin{itemize}[leftmargin=*]
  \item \textbf{Temperature lower bound}
        (Theorem~\ref{thm:absorption}):
        $\tau \ge C_\alpha \sigma_w \sqrt{\log p}$
        to avoid dead features.
  \item \textbf{Pruning time}
        (Theorem~\ref{thm:separation} +
        Proposition~\ref{prop:pruning}):
        $t^* = O(\tau^2 p / (c\kappa_{\min}))$.
  \item \textbf{Resulting time complexity}:
        $t^* = O(\sigma_w^2 p \log p / (c\kappa_{\min}))$.
\end{itemize}
The $p \log p$ scaling is a concrete, non-trivial prediction:
feature selection in $p$-dimensional soft stumps takes at least
$\Omega(p \log p)$ gradient-flow time with random initialization,
or $\Omega(p)$ with symmetric initialization.
This is the quantitative price of sparsity-inducing gating.

\subsection{Residual-Driven Specialization}
\label{sec:specialization}

The results above concern the \emph{early phase}, where the
mean-field interaction is weak ($f_\mu \approx 0$).
We now show that \emph{beyond} the early phase, the residual
coupling between stumps produces a qualitatively new phenomenon:
\textbf{dynamic rebalancing of gradient signal across features}.
This is a genuinely collective effect that does not occur without
ensemble interaction.

\begin{definition}[Effective signal]
  \label{def:eff-signal}
  For $j \in S$, the \emph{effective signal} at time $t$ is
  \begin{equation}
    \kappa_j^{\mathrm{eff}}(t)
      \coloneqq \E[r_j^{(\mu_t)}(x_j) \cdot x_j]
      = \kappa_j - \E\bigl[\E[f_{\mu_t}(\vx) \mid x_j]\cdot x_j\bigr].
    \label{eq:eff-signal}
  \end{equation}
\end{definition}

Since $f_{\mu_t}(\vx)$ decomposes into contributions from
each feature's stump group (by independence), we have
$\E[f_{\mu_t} \mid x_j] = f_{\mu_t}^{(j)}(x_j) + \mathrm{const}$,
where
$f_{\mu_t}^{(j)}(x_j) = \int_{\mathrm{feat}=j}
c\,\sigma_\beta(x_j - b)\,\dmu_t$
is the ensemble's output on feature $j$.
Thus:
\begin{equation}
  \kappa_j^{\mathrm{eff}}(t)
    = \kappa_j - \E[f_{\mu_t}^{(j)}(x_j)\cdot x_j].
  \label{eq:eff-signal-decomp}
\end{equation}

As the stumps on feature $j$ learn $g_j$,
$\E[f_{\mu_t}^{(j)} x_j]$ increases toward $\kappa_j$,
and $\kappa_j^{\mathrm{eff}}(t)$ decreases toward zero.
This is the \emph{residual feedback}: learning feature~$j$
\emph{weakens its own gradient signal}, redirecting gradient toward
less-learned features.

\begin{proposition}[Signal Crossing under Residual Coupling]
  \label{prop:crossing}
  Consider $k = 2$ features with
  $\kappa_1 > \kappa_2 > 0$.
  Suppose the mean-field dynamics is past the pruning phase
  (support is $S = \{1,2\}$) and the stump group on feature $1$
  makes progress in the sense that
  $t \mapsto \E[f_{\mu_t}^{(1)}(x_1)\,x_1]$ is strictly increasing.
  Then:
  \begin{enumerate}[label=(\alph*)]
    \item \textbf{Signal decay:}
          $\kappa_1^{\mathrm{eff}}(t)$ is strictly decreasing.
    \item \textbf{Signal crossing:}
          There exists $t_c > 0$ such that
          $\kappa_1^{\mathrm{eff}}(t_c) = \kappa_2^{\mathrm{eff}}(t_c)$,
          provided $\kappa_2^{\mathrm{eff}}$ has not yet fully
          decayed by time~$t_c$.
          For $t > t_c$,
          $\kappa_2^{\mathrm{eff}}(t) > \kappa_1^{\mathrm{eff}}(t)$:
          the initially weaker feature becomes the stronger signal.
    \item \textbf{Independent training fails to cross:}
          Under independent training
          (each stump minimizes $\E[(y - T(\vx;\,\vtheta))^2]$
          without residual sharing), the effective signal is
          $\kappa_j^{\mathrm{eff}} \equiv \kappa_j$ for all $t$,
          so $\kappa_1^{\mathrm{eff}} > \kappa_2^{\mathrm{eff}}$
          for all~$t$.
          No crossing occurs.
  \end{enumerate}
\end{proposition}

\begin{proof}
  Part (a) follows directly from~\eqref{eq:eff-signal-decomp}:
  $\kappa_1^{\mathrm{eff}}(t)
  = \kappa_1 - \E[f_{\mu_t}^{(1)} x_1]$,
  and $\E[f_{\mu_t}^{(1)} x_1]$ is increasing by assumption.

  Part (b): Initially,
  $\kappa_1^{\mathrm{eff}}(0) = \kappa_1 > \kappa_2
  = \kappa_2^{\mathrm{eff}}(0)$.
  If feature~$1$ is learned first (which it is, since
  $\kappa_1 > \kappa_2$), then $\kappa_1^{\mathrm{eff}}$
  decreases while $\kappa_2^{\mathrm{eff}}$ remains close
  to $\kappa_2$ (feature~$2$ is not yet being learned
  significantly).
  By the intermediate value theorem, there exists $t_c$ such
  that $\kappa_1^{\mathrm{eff}}(t_c) = \kappa_2^{\mathrm{eff}}(t_c)$.
  After $t_c$, the inequality reverses.

  Part (c): Without residual sharing, each stump independently
  targets $f^*(\vx)$.
  The gradient signal for feature $j$ in a single stump is
  $G_j = \E[f^*(\vx)\,\sigma_\beta'(\va^\top\vx-b)\,x_j]$,
  which is independent of other stumps' parameters.
  The effective signal is fixed at
  $\kappa_j^{\mathrm{eff}} = \kappa_j$ for all $t$.
\end{proof}

\begin{remark}[Why this requires mean-field theory]
  The signal crossing is a \emph{collective} phenomenon:
  each stump's gradient depends on the aggregate output $f_\mu$,
  which in turn depends on the distribution of all stumps.
  This circular dependence is precisely the mean-field interaction.
  Without it (independent training), there is no feedback loop
  and no rebalancing.
  %
  The practical consequence is that ensemble training with shared
  residuals---as in gradient boosting and mean-field dynamics---produces
  feature-level specialization that independent training cannot
  achieve.
\end{remark}

% ============================================================
\section{Proofs}
\label{sec:proofs}
% ============================================================

\subsection{Proof of Theorem~\ref{thm:separation}}

\paragraph{Setup.}
Write $\va = \entmax_\alpha(\vw/\tau)$ with
$a_j = 1/p + O(W/(\tau p))$ under Assumption~\ref{asm:hightemp}.
For each $j$, decompose $\va^\top\vx = a_j x_j + Z_j$ where
$Z_j = \sum_{l \neq j} a_l x_l \perp x_j$.

Define the conditional signal
$\phi_j(x_j) = \E_{\vx_{-j}}[r_\mu(\vx) \cdot
\sigma_\beta'(a_j x_j + Z_j - b)]$,
so that $G_j = \E_{x_j}[\phi_j(x_j) \cdot x_j]$.

\paragraph{Step 1: Taylor expansion.}
Since $a_j = O(1/p)$:
$\sigma_\beta'(a_j x_j + Z_j - b)
= \sigma_\beta'(Z_j - b) + \sigma_\beta''(Z_j - b)\,a_j x_j
+ O(x_j^2/p^2)$.

The mean-field predictor satisfies
$f_\mu(\vx) = f_\mu^{(0)}(\vx_{-j}) + h_j(\vx_{-j})\,x_j
+ O(x_j^2/p^2)$,
where $h_j = \int c' a_j' \sigma_\beta'(Z_j' - b')\dmu = O(C/(4\beta p))$
pointwise, and the remainder bound
$|O(x_j^2/p^2)| \le Cx_j^2/(4\beta^2 p^2)$
holds uniformly over $\mu$ supported on $\Theta$.

\paragraph{Step 2: Part (b) ($l \notin S$).}
$\phi_l(x_l) = \phi_l^{(0)} + \phi_l'(0)\,x_l + O(x_l^2/p^2)$
where $\phi_l^{(0)} = \E_{\vx_{-l}}[(f^*(x_S) - f_\mu^{(0)})
\sigma_\beta'(Z_l - b)]$
is independent of $x_l$.
Since $\E[x_l] = 0$:
$G_l = \phi_l^{(0)} \cdot 0 + \phi_l'(0) \cdot \E[x_l^2]
+ O(\E[|x_l|^3]/p^2) = \phi_l'(0) + O(1/p^2)$.

Now $\phi_l'(0) = a_l\,\E[(f^*{-}f_\mu^{(0)})\sigma_\beta''(Z_l{-}b)]
- \E[h_l\,\sigma_\beta'(Z_l{-}b)]$.
Both terms are $O(1/p)$:
$|a_l| \le 1/p + O(W/(\tau p))$ and $|h_l| \le C/(4\beta p)$.
So $|G_l| \le C_2/p$.
This proves part (b).

\paragraph{Step 3: Part (a) ($j \in S$).}
The residual contains $g_j(x_j)$:
$\phi_j(x_j)
= g_j(x_j)\,\E_{Z_j}[\sigma_\beta'(Z_j{-}b)]
+ (\text{terms independent of } x_j) + O(x_j/p)$.

Hence $G_j = \E[g_j(x_j)x_j]\,\E[\sigma_\beta'(Z_j{-}b)]
- \E[\E[f_\mu \mid x_j]\,x_j]\,\E[\sigma_\beta'(Z_j{-}b)]
+ O(1/p)$.

During the early phase, $\|f_\mu\|_{L^2} \le \kappa_{\min}/2$, so
by Jensen's inequality
$|\E[\E[f_\mu \mid x_j]\,x_j]| \le \|f_\mu\|_{L^2}
\le \kappa_{\min}/2$.
Thus $\E[r_j^{(\mu)}(x_j)\,x_j] \ge |\kappa_j| - \kappa_{\min}/2
\ge \kappa_{\min}/2$.
This proves part (a).

\paragraph{Step 4: Parts (c) and (d).}
$\dot{w}_j = \frac{c}{\tau}(\frac{G_j}{p}
- \frac{\sum_l G_l}{p^2}) + O(1/p^2)
= \frac{c}{\tau p}\,G_j + O(1/p^2)$
for $j \in S$ (since $\sum_l G_l = O(k)$ gives
$\sum G_l/p^2 = O(1/p^2)$).
For $l \notin S$: $G_l = O(1/p)$, so $\dot{w}_l = O(1/p^2)$.
The gap dynamics follows. \qed

% ============================================================
\section{Discussion}
\label{sec:discussion}
% ============================================================

\paragraph{What this theory uncovers.}
The qualitative advice ``explore broadly, then sparsify'' is
common wisdom.
The contribution of this paper is not that advice, but the
identification of precise dynamical mechanisms that underlie it:
(i)~absorbing dead-feature states that give a
\emph{lower bound on temperature}
(Theorem~\ref{thm:absorption});
(ii)~a $O(1/(\tau p))$ gap-growth bottleneck that gives an
\emph{upper bound on selection speed}
(Theorem~\ref{thm:separation});
and (iii)~residual-driven signal crossing that shows
why ensemble coupling produces feature-level specialization
(Proposition~\ref{prop:crossing}).
The interplay of (i) and (ii) yields a concrete
time--temperature tradeoff ($t^* = O(\tau^2 p)$);
result~(iii) shows that this tradeoff is not merely a constraint
but part of a richer dynamics in which the ensemble self-organizes
across features.

\paragraph{Implications for model design.}
The $\Omega(p)$ lower bound on feature-selection time
(with symmetric initialization) reveals a structural disadvantage
of gradient-trained soft trees relative to greedy (hard) trees:
greedy splitting evaluates each feature independently and does
not pay the $1/p$ dilution cost.
This may partially explain the continued practical dominance of
hard tree ensembles on high-dimensional tabular
data~\citep{grinsztajn2022tree}.
On the other hand, Proposition~\ref{prop:crossing} shows a
potential advantage of soft-tree ensembles: the residual coupling
creates feature-level specialization that could lead to
more balanced coverage of multi-feature targets than
independent stump training.

\paragraph{Limitations.}
Our analysis of feature discovery (Theorems~\ref{thm:absorption}
and~\ref{thm:separation}) is restricted to the early phase.
Proposition~\ref{prop:crossing} goes beyond the early phase
but assumes that in-support features make progress; a fully
coupled gap--risk analysis remains open.
Throughout, we assume depth-one trees, additive targets with
independent covariates, and first-order signal ($\kappa_j \neq 0$).
Extending to deeper trees, feature interactions, and higher
information exponents are important open directions.

% ============================================================
\begin{thebibliography}{99}

\bibitem[Blondel et~al.(2020)]{blondel2020learning}
M.~Blondel, A.~F.~T. Martins, and V.~Niculae.
\newblock Learning with {Fenchel-Young} losses.
\newblock \emph{JMLR}, 21(35):1--69, 2020.

\bibitem[Chizat and Bach(2018)]{chizat2018global}
L.~Chizat and F.~Bach.
\newblock On the global convergence of gradient descent for
  over-parameterized models using optimal transport.
\newblock In \emph{NeurIPS}, 2018.

\bibitem[Frosst and Hinton(2017)]{frosst2017distilling}
N.~Frosst and G.~Hinton.
\newblock Distilling a neural network into a soft decision tree.
\newblock In \emph{CEx Workshop at AI*IA}, 2017.

\bibitem[Grinsztajn et~al.(2022)]{grinsztajn2022tree}
L.~Grinsztajn, E.~Oyallon, and G.~Varoquaux.
\newblock Why do tree-based models still outperform deep learning on
  typical tabular data?
\newblock In \emph{NeurIPS}, 2022.

\bibitem[Hazimeh et~al.(2020)]{hazimeh2020tree}
H.~Hazimeh, N.~Ponomareva, P.~Mol, Z.~Tan, and R.~Mazumder.
\newblock The tree ensemble layer: Differentiability meets conditional
  computation.
\newblock In \emph{ICML}, 2020.

\bibitem[Mei et~al.(2018)]{mei2018mean}
S.~Mei, A.~Montanari, and P.-M. Nguyen.
\newblock A mean field view of the landscape of two-layer neural
  networks.
\newblock \emph{PNAS}, 115(33):E7665--E7671, 2018.

\bibitem[Peters et~al.(2019)]{peters2019sparse}
B.~Peters, V.~Niculae, and A.~F.~T. Martins.
\newblock Sparse sequence-to-sequence models.
\newblock In \emph{ACL}, 2019.

\bibitem[Popov et~al.(2020)]{popov2020neural}
S.~Popov, S.~Morozov, and A.~Babenko.
\newblock Neural oblivious decision ensembles for deep learning on
  tabular data.
\newblock In \emph{ICLR}, 2020.

\end{thebibliography}

\end{document}
